\documentclass[12pt]{article}
\usepackage{amsmath,amssymb}

\begin{document}
\section*{{2025 AMC 12B Problems/Problem 8}}
Source: AoPS Wiki

\subsection*{Solution}
Because every coefficient of our polynomial $f(x)$ is an integer, the conjugate of $4+\sqrt{5}$ , $4-\sqrt{5}$ is also a root. Let the third and final root be $r_3$ . Thanks to Vieta's Formulas, we know that the sum of the roots is \[(r_3) + (4+\sqrt{5}) + (4 - \sqrt{5}) = 5\] \[r_3 = -3.\] We now calculate $f(1) = (1+3)(1-4-\sqrt{5})(1-4+\sqrt{5})$ . $(1-4-\sqrt{3})(1-4+\sqrt{3}) = (-3-\sqrt{5})(-3+\sqrt{5}) = 9 - 5 = 4.$ So $f(1) = 4 \cdot 4 = 16.$ Hence $f(1) = 16 = 1-5 + a + b$ . Therefore, $a + b = \boxed{20}$ . ~lprado (minor edits by singularity1)

\end{document}
